\chapter{Related Work}
\label{chap:related-work}

The intrusion detection landscape is dense. To position \nidsname{}
honestly I divided the comparison into three groups: production-grade
open-source NIDS engines, commercial SIEM platforms and the academic
machine-learning IDS literature. Each section ends with a short note
on how my project differs.

\section{Snort}
\label{sec:rw-snort}

I'll start with the tool everyone benchmarks against.
\textbf{Snort}~\cite{roesch1999snort} arrived in 1998 from Martin
Roesch as the open-source answer to expensive commercial NIDS
appliances of the late 1990s. Almost three decades later it is
still under active development; the current C++ rewrite (Snort 3,
2021) replaced the original C engine that Roesch wrote, but the
detection philosophy hasn't shifted: rules in, alerts out, no UI
between them.

What people actually buy into when they pick Snort is the rules
ecosystem rather than the engine. Cisco's Talos team (Cisco bought
Sourcefire, the company behind Snort, back in 2013) ships a daily
Subscriber rule set for paying customers and a weekly Community
one for free users. Proofpoint runs Emerging Threats, a
third-party project that distributes thousands of additional
rules at no charge. A real-world Snort deployment is essentially
a subscription to detection logic written by other people --- the
engine's biggest strength on day one, and the recurring tax on
day two when the rules need pruning, tuning and retiring.

On the engineering side Snort was built for line-rate inspection.
With the right NIC and config it'll happily process multiple
gigabits per second of inline traffic. The cost of that
throughput is an interface aimed squarely at terminal-comfortable
operators. There's no analyst console in the box. Snort writes
unified2 binary logs that some downstream pipeline
(Barnyard2 typically, then a SIEM, then a dashboard) has to chew
through before anyone sees an alert. Operators who want a UI pick
between SecurityOnion (an entire Linux distribution that bundles
Snort with half a dozen other tools), the commercial Snort Cloud
offering from Cisco, or some homegrown stack.

The shape of \nidsname{}'s signature engine borrows from Snort,
but two things differ. The match specification is JSON instead of
a custom DSL, and it's editable from the UI rather than a text
file in \file{/etc/snort/}. And the analyst console ships
in the same binary as the detection engine instead of being
somebody else's problem.

\section{Suricata}
\label{sec:rw-suricata}

The number-two engine in the open-source NIDS world is
\textbf{Suricata}~\cite{suricata}. The Open Information Security
Foundation (OISF) started it in 2009 with a clean-slate
multi-threaded rewrite of the same ideas Snort had pioneered a
decade earlier. The trick the OISF pulled was making Suricata
read Snort rules (with a few extensions of its own), so the rule
ecosystem is shared --- if you write detection for one engine you
get the other almost for free.

Three things make Suricata interesting from an engineering angle.
First, threads: it scales close to linearly with CPU cores out of
the box, where Snort needed serious tuning to do the same.
Second, protocol parsing: HTTP, TLS, SSH, SMB and DNS are decoded
before the matching happens, so a rule can fire on parsed
fields (\code{tls.sni}, \code{http.uri}) rather than raw
payload bytes. Third, the Eve log format --- one JSON object per
event, structured, easy for downstream SIEMs to ingest. Compared
to Snort's unified2 binary blob, that's a real quality-of-life
upgrade.

Like Snort, though, Suricata is a daemon. OISF publishes Scirius
on the side for rule management and a separate community project
EveBox for alert browsing, but the engine itself ships without an
analyst-facing UI.

I borrowed Suricata's structured-logging idea --- JSON in,
JSON out, no binary intermediates --- but stopped short of the
protocol-aware matching it's famous for. The rule vocabulary in
\nidsname{} works at the TCP/UDP/ICMP level on headers and
payloads, which covers the validation scenarios (port scans,
brute-force attempts, simple floods, content-based exploits).
What it can't do is the parser-level matching that's needed for
HTTP request smuggling, TLS-fingerprint analysis or DNS
tunnelling --- those would require pulling in a stateful
protocol parser layer, which is its own project.

\section{Zeek}
\label{sec:rw-zeek}

\textbf{Zeek}~\cite{paxson1999bro} (the project everyone still
calls Bro, two decades after Vern Paxson started it under that
name) belongs to a different category altogether. Snort and
Suricata are pattern matchers. Zeek is a programming environment
for the network. Operators write scripts in Zeek's own scripting
language and hang them on events the engine emits ---
\code{connection\_established}, \code{http\_request},
\code{dns\_request}, and a long list of others --- and the
scripts decide what to log. Out of the box the engine doesn't
produce ``alerts'' in the Snort sense. It produces structured
log streams (\file{conn.log}, \file{http.log}, \file{dns.log},
and dozens more) and leaves the detection logic to whatever the
operator scripts or to a downstream SIEM.

That design is why threat-hunters love it. If your job is asking
exploratory questions of the traffic record --- ``which hosts
talked to this domain in the last 30 days, on what ports, with
what user-agent strings?'' --- Zeek is the right answer. Its
logs are also the de-facto input format for academic ML-IDS
research. CICIDS2017 and IOT-23, the datasets every recent
paper benchmarks against, ship Zeek log bundles alongside the
raw PCAP.

\nidsname{} doesn't try to compete in this space. The audience
this project targets isn't the researcher who wants to script
arbitrary queries over weeks of traffic. It's the operator who
wants the console to say ``alert fired, here it is, here are
the packets that produced it, what do you want to do''. The two
tools could live next to each other in a real deployment ---
Zeek doing the forensic logging, \nidsname{} doing the
real-time alerting --- without overlap.

\section{Wazuh}
\label{sec:rw-wazuh}

\textbf{Wazuh}~\cite{wazuh} sits a tier up the stack. It
grew out of OSSEC --- the long-running host-based IDS --- and
added a central manager, a Kibana-based console and a pile of
integrations with external network engines. The detection that
Wazuh itself does is primarily host-side: log analysis, file
integrity monitoring, vulnerability assessment, audit-rule
matching. For network-layer detection it relies on a sidecar
that ingests alerts from Snort or Suricata --- Wazuh aggregates
and dashboards them, but it doesn't generate them.

What Wazuh ships is a full operations console: rule management,
incident queues, agent administration, configurable dashboards
on top of Elasticsearch (or its OpenSearch fork). All of it is
production-grade. The price you pay is infrastructure footprint.
A minimal Wazuh single-node deployment wants several gigabytes
of RAM just to keep the search index responsive, and the
recommended cluster setup runs three to five nodes.

\nidsname{} differs from Wazuh in the same way a single-binary
SQLite app differs from a clustered PostgreSQL deployment: by an
order of magnitude in setup complexity. Wazuh would be the right
choice for an enterprise. \nidsname{} is the right choice for a
laptop demo or a small-network deployment where the SIEM is allowed
to be one process and one database file.

\section{Splunk Enterprise Security}
\label{sec:rw-splunk}

The benchmark in the commercial SIEM market is
\textbf{Splunk Enterprise Security}. Architecturally it works
much like Wazuh --- agents on hosts feed events into a central
indexer; analysts query that indexer from a web console;
saved correlation searches roll up matched events into
``notable events'' that feed an incident queue. What separates
Splunk is the polish of the analyst experience. When IBM QRadar
or Microsoft Sentinel get reviewed in trade press, the comparison
is almost always against Splunk.

The cost of that polish is the other reason Splunk gets talked
about. Licensing is volume-priced --- you pay per gigabyte of
data ingested per day --- and even a single-instance ES
deployment lands in five-figure annual budgets. There is a free
Splunk Free tier, but it caps ingest at 500\,MB per day and
doesn't include the Enterprise Security app at all.

So \nidsname{} is not in the same market. It sits at the
opposite end of the spectrum: free, self-hosted, single binary,
for environments where the budget is zero and the deployment
needs to take minutes instead of weeks.

\section{Academic Machine-Learning IDS}
\label{sec:rw-ml-academia}

The published machine-learning IDS literature is, as discussed in
Section~\ref{sec:bg-ml}, vast. The bulk of it consists of papers
that propose a particular classifier architecture (a deep neural
network, a gradient-boosted ensemble, an autoencoder, a transformer)
and evaluate it on a standardised dataset. The two most-cited
datasets in recent work are
\textbf{NSL-KDD}~\cite{tavallaee2009detailed} (a cleaned version of
the 1999 KDD Cup dataset) and
\textbf{CICIDS2017}~\cite{sharafaldin2018cicids} (a more recent
dataset from the Canadian Institute for Cybersecurity).

The Sommer and Paxson critique~\cite{sommer2010outside} is now
fifteen years old but remains relevant. They argue, persuasively,
that the closed-world assumption of standard machine-learning
evaluation does not transfer to operational network security. A
classifier that scores 0.99 F1 on NSL-KDD is not necessarily a
useful production tool, because the distribution of normal traffic
in NSL-KDD is not the distribution of normal traffic on any real
network. Sommer and Paxson recommend three concrete safeguards:
prefer anomaly detection over misuse detection for the ML
component, give the operator a clear explanation of why each alert
fired, and keep the ML engine in a supporting role behind a
signature engine.

\nidsname{} explicitly adopts all three recommendations. Isolation
Forest is offered as a first-class option for operators who lack
labelled attack data; every alert is tagged with its detection
method; and the hybrid category is the high-confidence one,
implicitly demoting the ml-only category to a ``second opinion''
role.

\section{Positioning of This Work}
\label{sec:rw-positioning}

The position \nidsname{} occupies is most easily summarised in a
feature matrix; Table~\ref{tab:positioning} reproduces the
comparison from the presentation deck.

\begin{table}[h]
\centering
\small
\caption{Feature comparison of \nidsname{} against related tools.
``$\checkmark$'' indicates the feature is present out of the box,
``$\sim$'' indicates partial or plugin-based support, ``--''
indicates the feature is outside the tool's scope, and ``$\times$''
indicates the feature is absent.}
\label{tab:positioning}
\begin{tabular}{l c c c c c c}
\toprule
                              & Snort & Suricata & Zeek & Wazuh & Splunk ES & \textbf{Ovidius NIDS} \\
\midrule
Signature engine              & $\checkmark$ & $\checkmark$ & $\sim$ & $\sim$  & $\sim$  & $\checkmark$ \\
ML engine                     & $\times$     & $\times$     & $\times$ & $\sim$  & $\checkmark$ & $\checkmark$ \\
Hybrid label                  & --           & --           & --     & $\times$ & $\times$ & $\checkmark$ \\
Built-in analyst UI           & $\times$     & $\times$     & $\times$ & $\checkmark$ & $\checkmark$ & $\checkmark$ \\
Per-alert PCAP                & $\times$     & $\times$     & $\times$ & $\times$ & $\times$ & $\checkmark$ \\
Real-time stream UI           & $\times$     & $\times$     & $\times$ & $\sim$   & $\checkmark$ & $\checkmark$ \\
Editable rules from UI        & $\times$     & $\times$     & $\times$ & $\checkmark$ & $\checkmark$ & $\checkmark$ \\
One-click deploy (Win/Mac/Linux)& $\times$   & $\times$     & $\times$ & $\times$ & $\times$ & $\checkmark$ \\
Self-hostable                 & $\checkmark$ & $\checkmark$ & $\checkmark$ & $\checkmark$ & $\sim$ & $\checkmark$ \\
Free / open-source            & $\checkmark$ & $\checkmark$ & $\checkmark$ & $\checkmark$ & $\times$ & $\checkmark$ \\
\bottomrule
\end{tabular}
\end{table}

The differentiators that matter for a thesis-scale project, and that
no other tool in the table provides simultaneously, are:

\begin{itemize}
  \item Detection, ML, forensics and UI in a single binary.
  \item Per-alert packet capture as a core feature rather than a
    bolt-on integration.
  \item One-click deployment with no infrastructure dependencies.
  \item Hybrid alert labelling that distinguishes signature, ML and
    cooperative detection visually.
\end{itemize}

The trade-off is scale. \nidsname{} is not built to inspect
multi-gigabit traffic on a forty-core appliance; it is built to
run on a laptop and demonstrate the cooperation of the three
engines on real packets. That is the niche, and the comparison
table makes it explicit.
